\BourbakiExerciseGroup

\begin{enumerate}
\def\labelenumi{\arabic{enumi})}
\item
  Let E be an extension of a field K and x and y two distinct roots in E of the same irreducible polynomial of K{[}X{]}. Show that the extensions \(\mathbf{K}(x)\) and \(\mathbf{K}(y)\) are not linearly disjoint over K (use Th. 2). * Give an example where K = Q and \(\mathbf{K}(x) \cap \mathbf{K}(y) = \mathbf{K}\) (cf.~V, p.~161, Ex. 2). *
\item
  Let \((\mathbf{E}_i)_{i\in I}\) be a family of extensions of a field \(\mathbf{K}\) , contained in an extension \(\mathbf{L}\) of \(\mathbf{K}\) . If \(\mathbf{F}_i\) is the relative algebraic closure of \(\mathbf{K}\) in \(\mathbf{E}_i\) , show that the relative algebraic closure of \(\mathbf{K}\) in \(\mathbf{E} = \bigcap_{i\in I}\mathbf{E}_i\) is \(\mathbf{F} = \bigcap_{i\in I}\mathbf{F}_i\) .
\item
  Show that every algebraic extension E of a field K is equipotent to a subset of \(K \times N\) (consider the mapping which associates with each element of E its minimal polynomial over K). In particular, every algebraic extension of a finite field is countable and every algebraic extension of an infinite field K is equipotent to K. * Deduce that in the field R of real numbers there exist numbers that are transcendental over the prime subfield Q and that the set of all these numbers has the power of the continuum. *
\item
  Let F be a transcendental extension of a field E. Show that
\end{enumerate}

\[
[ \mathrm{F}: \mathrm{E} ] \geqslant \operatorname * {S u p} (\operatorname{Card} (\mathrm{E}), \operatorname{Card} (\mathbf {N}))
\]

with equality if the extension F can be generated by a countable family of elements. (Use Ex. 7 of IV, p.~89.)

\begin{enumerate}
\def\labelenumi{\arabic{enumi})}
\setcounter{enumi}{4}
\tightlist
\item
  Let A be a commutative algebra and let S be a subset of A generating A as K-algebra. Suppose that \(\operatorname{Card}(S) < \operatorname{Card}(K)\) . Show that if M is a maximal ideal of A, then the field A/M is an algebraic extension of K. (Use Ex. 15 of V, p.174 when \(\operatorname{Card}(K) \leqslant \operatorname{Card}(N)\) and Ex. 4 when
\end{enumerate}

\[
\operatorname{Card} (\mathbf {K}) > \operatorname{Card} (\mathbf {N}).
\]

* Deduce that \(\mathbf{A} / \mathfrak{M} = \mathbf{K}\) when \(\mathbf{K}\) is algebraically closed. *

* 6) Let C be an uncountable algebraically closed field and \(\mathbf{P} = \mathbf{C}[(\mathbf{X}_i)_{i \in I}]\) a polynomial algebra over C. Given two families \((f_\alpha)_{\alpha \in A}, (g_\beta)_{\beta \in B}\) of elements of P suppose that: a) B and I are countable.

\begin{enumerate}
\def\labelenumi{\alph{enumi})}
\setcounter{enumi}{1}
\tightlist
\item
  For every finite subset \(\mathbf{A}'\) of \(\mathbf{A}\) and every finite subset \(\mathbf{B}'\) of \(\mathbf{B}\) there exists \((x_i)_{i\in I} \in C^I\) such that
\end{enumerate}

\[
f _ {\alpha} (x _ {i}) = 0 \quad \text { and } \quad g _ {\beta} (x _ {i}) \neq 0 \quad \text { for   all } \quad \alpha \in \mathbf {A} ^ {\prime}, \beta \in \mathbf {B} ^ {\prime}.
\]

Show that then there exists \((x_{i})_{i\in I}\in \mathbf{C}^{I}\) such that

\[
f _ {\alpha} (x _ {i}) = 0 \quad \text { and } \quad g _ {\beta} (x _ {i}) \neq 0 \quad \text { for   all } \quad \alpha \in A, \beta \in B.
\]

(Let \((\mathbf{Y}_{\beta})_{\beta \in \mathbb{B}}\) be a family of indeterminates indexed by \(\mathbf{B}\) and let \(\mathbf{Q}\) be the polynomial algebra \(\mathbf{C}[(\mathbf{X}_i)_{i\in I}, (\mathbf{Y}_{\beta})_{\beta \in \mathbb{B}}]\) . Let \(\Lambda\) be the quotient of \(\mathbf{Q}\) by the ideal generated by the \(f_{\alpha}\) and the \(1 - g_{\alpha}\mathbf{Y}_{\alpha}\) . Show that \(\Lambda \neq 0\) thanks to \(a)\) and \(b)\) and apply Ex. 5 to a maximal ideal of \(\Lambda\) .) *

\begin{enumerate}
\def\labelenumi{\arabic{enumi})}
\setcounter{enumi}{6}
\tightlist
\item
  Let \(L\) be a field, \(A\) a subring of \(L\) and \(K \subset L\) the field of fractions of \(A\) .
\end{enumerate}

\begin{enumerate}
\def\labelenumi{\alph{enumi})}
\item
  Show that if L is a finitely generated A-module (II, p.~206) then necessarily A = K (by considering a supplementary subspace of K in L, qua vector space over K, show that K is a finitely generated A-module, hence necessarily of the form \(s^{-1}A\) , where \(s \in A\) ; finally show that s must be invertible in A).
\item
  Show that if there exist a finite number of elements \(x_{j}\) ( \(1 \leqslant j \leqslant n\) ) of L, algebraic over K, such that \(L = A[x_{1}, \ldots, x_{n}]\) , then there exists an element \(b \neq 0\) in A such that \(K = A[b^{-1}]\) (show that there exists \(b \neq 0\) in A such that L is a finitely generated \(A[b^{-1}]\) -module). Show that either b is invertible (and K = A) or every non-zero ideal of A contains a principal ideal \(Ab^{m}\) . Converse. The ring \(k[[X]]\) of formal power series over a field k is a ring of this type.
\end{enumerate}

* 8) If K is a field, show that K is algebraically closed in the field K(X) of rational fractions in one indeterminate over K (use the fact that K{[}X{]} is a principal ideal domain). *

\begin{enumerate}
\def\labelenumi{\arabic{enumi})}
\setcounter{enumi}{8}
\tightlist
\item
  Let \(\mathbf{K}\) be a field, \(\mathbf{E}\) an extension of \(\mathbf{K}\) and \(\mathbf{S}\) a subset of \(\mathbf{E} - \mathbf{K}\) .
\end{enumerate}

\begin{enumerate}
\def\labelenumi{\alph{enumi})}
\item
  Show that the set of extensions \(\mathbf{L} \subset \mathbf{E}\) of \(\mathbf{K}\) such that \(\mathbf{L} \cap \mathbf{S} = \varnothing\) , ordered by inclusion, has a maximal element \(\mathbf{M}\) .
\item
  Show that if S is finite, E is an algebraic extension of M.
\end{enumerate}

* 10) Let K be a field and P a polynomial in the ring K{[}X, Y{]} in two indeterminates. Suppose that there exists a polynomial H ∈ K{[}X{]} and two polynomials Q, R ∈ K{[}X, Y{]} such that H(X) P(X, Y) = Q(X, Y) R(X, Y) and that the coefficients of Q(X, Y), qua polynomial in Y, are polynomials in K{[}X{]} without non-constant common factor. Show that all the coefficients of R(X, Y), qua polynomial in Y, are divisible by H(X). (Consider H, P, Q, R as polynomials in X over the field K(Y) and decompose them into factors of the first degree in Ω{[}X{]} where Ω is an algebraic closure of K(Y); observe that if H(X) and Q(X, Y) had a factor of the first degree in common, this factor would divide all the coefficients of Q(X, Y), qua polynomial in Y, and use the fact (IV, p.~12) that if K' is an extension of K, the GCD in K'{[}X{]} of two polynomials of K{[}X{]} belongs to K{[}X{]}.) *

\begin{enumerate}
\def\labelenumi{\arabic{enumi})}
\setcounter{enumi}{10}
\tightlist
\item
  Let \(\mathbf{E}\) be an extension of a field \(\mathbf{K}, x \in \mathbf{E}\) a transcendental element over \(\mathbf{K}\) , so that \(\mathbf{K}(x)\) is isomorphic to the field \(\mathbf{K}(\mathbf{T})\) of rational fractions in an indeterminate over \(\mathbf{K}\) . Thus every \(y \in \mathbf{K}(x)\) may be written \(y = g(x) / h(x)\) , where \(g\) and \(h\) are two polynomials of \(\mathbf{K}[\mathbf{X}]\) that are relatively prime (IV, p.~12), determined up to a non-zero common factor in \(\mathbf{K}\) ; the larger of the degrees of \(g\) and \(h\) is called the height of \(y\) with respect to \(x\) .
\end{enumerate}

\begin{enumerate}
\def\labelenumi{\alph{enumi})}
\tightlist
\item
  Show that in the polynomial ring \(\mathbf{K}(y)[\mathbf{X}]\) the polynomial \(g(\mathbf{X}) - yh(\mathbf{X})\) is irreducible if \(y \notin \mathbf{K}\) . Deduce that if \(y\) has height \(n > 0\) with respect to \(x\) , then \(\mathbf{K}(x)\) is an algebraic extension of degree \(n\) of \(\mathbf{K}(y)\) .
\end{enumerate}

Let \(\mathrm{P}(\mathbf{X}) = \sum_{i=0}^{n} y_i \mathbf{X}^i\) be the minimal polynomial of \(x\) over \(\mathrm{K}(y)\) . Show that for

\(0 \leqslant i \leqslant n\) we have either \(y_i \in \mathbf{K}\) or \(\mathbf{K}(y_i) = \mathbf{K}(y)\) .

\begin{enumerate}
\def\labelenumi{\alph{enumi})}
\setcounter{enumi}{1}
\item
  Deduce from a) that every \(y \in \mathbf{K}(x)\) such that \(\mathbf{K}(y) = \mathbf{K}(x)\) is of the form \((ax + b)/(cx + d)\) where a, b, c, d are elements of K such that \(ad - bc \neq 0\) ; converse. Find all K-automorphisms of \(\mathbf{K}(x)\) .
\item
  Show that if \(y \in \mathbf{K}(x)\) is of height \(n\) with respect to \(x\) , and \(z \in \mathbf{K}(y)\) of height \(m\) with respect to \(y\) , then \(z\) is of height \(mn\) with respect to \(x\) .
\end{enumerate}

* \(d\) ) Let F be an extension of K such that \(\mathbf{K} \subset \mathbf{F} \subset \mathbf{K}(x)\) and \(\mathbf{F} \neq \mathbf{K}\) ; let y be an element of F whose height \(m\) with respect to \(x\) has the least possible value; show that \(\mathbf{F} = \mathbf{K}(y)\) («Lüroth's theorem »). (Let P be the minimal polynomial of \(x\) in F{[}T{]}; show that \(\mathrm{P(T)} = \mathrm{Q(T,x)}/\mathrm{R}(x)\) , where Q is a polynomial of K{[}T,X{]} of degree \(\geqslant m\) in X and is not divisible by any non-constant polynomial of K{[}X{]}, and R is a polynomial of K{[}X{]}. If \(y = g(x)/h(x)\) , where \(g\) and \(h\) are relatively prime in K{[}X{]}, observe that by \(a\) ) the polynomial \(g(T)h(X) - g(X)h(T)\) is not divisible by any non-constant polynomial of K{[}T{]} or of K{[}X{]}; deduce that necessarily

\[
g (T) h (X) - g (X) h (T) = c. Q (T, X), \quad \text { where } c \in K,
\]

using Ex. 10.) *

* 12) a) With the notation and hypotheses as in Ex. 11, show that for an extension F of K such that \(\mathbf{K} \subset \mathbf{F} \subset \mathbf{K}(x)\) to be such that \(\mathbf{F} \cap \mathbf{K}[x] \neq \mathbf{K}\) , it is necessary and sufficient that \(\mathbf{F} = \mathbf{K}(y)\) where \(y = g(x)\) for a non-constant polynomial \(g \in \mathbf{K}[\mathrm{T}]\) ; then we have \(\mathbf{K}(y) \cap \mathbf{K}[x] = \mathbf{K}[y]\) . (Write \(\mathbf{F} = \mathbf{K}(y)\) with \(y = g(x)/h(x)\) , \(g\) and \(h\) being relatively prime polynomials of \(\mathbf{K}[\mathrm{T}]\) . If there exists a non-constant polynomial \(\mathbf{P} \in \mathbf{K}[\mathrm{T}]\) such that \(\mathbf{P}(x) \in \mathbf{F}\) , we can write \(\mathbf{P}(x) = \mathbf{Q}(y)/\mathbf{R}(y)\) , where \(\mathbf{Q}\) and \(\mathbf{R}\) are two relatively prime monic polynomials of \(\mathbf{K}[\mathrm{T}]\) , not both constant. Distinguish two cases, according as \(\mathbf{R} = 1\) or \(\mathbf{R}\) is non-constant; in the second case, decompose \(\mathbf{Q}\) and \(\mathbf{R}\) into factors of the first degree in an algebraic closure \(\Omega\) of \(\mathbf{K}\) and observe that if \(a, b\) are distinct elements of \(\Omega\) , then \(g - ah\) and \(g - bh\) are relatively prime in \(\Omega[\mathrm{T}]\) .)

\begin{enumerate}
\def\labelenumi{\alph{enumi})}
\setcounter{enumi}{1}
\item
  If \(\mathbf{F} = \mathbf{K}(y)\) with \(y = g(x)\) , where \(g \in \mathbf{K}[T]\) is non-constant, we may suppose that \(g(T) = Tg_1(T)\) , where \(g_1 \in \mathbf{K}[T]\) . Show that if \(\mathbf{K}(xg(x)) \cap \mathbf{K}[y] \neq \mathbf{K}\) , we necessarily have \(g_1(T) = aT^n\) with \(a \in \mathbf{K}\) and \(n > 0\) . (Observe that by hypothesis and \(a\) ) there exist two non-constant polynomials \(P, Q\) in \(K[T]\) such that \(P(xg(x)) = Q(g(x)) \notin K\) ; deduce that there are two elements \(a, b\) in \(K\) and an integer \(m\) such that \(g(T)\) divides \(a - bT^m\) .) c) Let \(y, z\) be two elements of \(K[x]\) such that \(K(y) \cap K(z) \neq K\) ; show that then \(K[y] \cap K[z] \neq K\) .
\item
  Find \(y \in \mathbf{K}(x)\) such that \(\mathbf{K}(y) = \mathbf{K}(x)\) but \(\mathbf{K}[y] \cap \mathbf{K}[x] = \mathbf{K}\) . *
\end{enumerate}

* 13) Let K be a field, \(P(T, X) = a_0(T)X^n + a_1(T)X^{n-1} + \cdots + a_n(T)\) a polynomial in K{[}T, X{]}; suppose that: \(1^\circ\) the polynomial \(a_0\) is not divisible by T; \(2^\circ\) each of the polynomials \(a_1, \ldots, a_n\) is divisible by T; \(3^\circ\) the polynomial \(a_n\) is not divisible by \(T^2\) . Show that under these conditions P, as polynomial of K(T){[}X{]}, is irreducible (argue by contradiction, using Ex. 10). *

* 14: With the notation as in Ex. 11, show that if we take \(g(X) = X\) and \(h(X) = X^n + X + 1 \quad (n \geqslant 2)\) , the polynomial \((g(X)h(x) - h(X)g(x)) / (X - x)\) of \(K(x)[X]\) is irreducible (use Ex. 13). *

* 15) Give an example of a field K containing two subfields \(\mathbf{K}_1\) , \(\mathbf{K}_2\) such that K is algebraic over \(\mathbf{K}_1\) and over \(\mathbf{K}_2\) but not over \(\mathbf{K}_1 \cap \mathbf{K}_2\) . (Take \(\mathbf{K} = \mathbf{Q}(T)\) , \(\mathbf{K}_1 = \mathbf{Q}(T^2)\) , and \(\mathbf{K}_2 = \mathbf{Q}(T^2 - T)\) . Verify that \([\mathbf{K} : \mathbf{K}_1] = 2\) , \([\mathbf{K} : \mathbf{K}_2] = 2\) and \(\mathbf{K}_1 \cap \mathbf{K}_2 = \mathbf{Q}\) . To do this note that an element \(f(T)\) of \(\mathbf{K}_1 \cap \mathbf{K}_2\) satisfies the conditions \(f(-T) = f(T)\) and \(f(1 - T) = f(T)\) .)

